\color{unidarkgreen}{\subsection{Комбинированный пусковой орган напряжения (КПОН)}}
\color{black}

\begin{enumerate}[label=\arabic{section}.\arabic{subsection}.\arabic*,  align=left, labelsep=4pt, leftmargin=0pt, itemindent=57pt]

\item
Комбинированный пусковой орган по напряжению первой системы шин (КПОН 1СШ) и второй системы шин
(КПОН 2СШ) используется для повышения чувствительности ступеней МТЗ.  

\begin{figure}[H]
\centering
\includegraphics[width=0.7\textwidth,height=0.7\textheight,keepaspectratio]{\fbpath/661.0371 КПОН/img/KPON.pdf}
\captionof{figure}{Функциональная схема КПОН}\label{pic_kpon:KPON}
\end{figure}

\item
Внутренний комбинированный пусковой орган напряжения, контролирует минимальное линейное
напряжение и максимальное напряжение обратной последовательности от ТН, установленных на первой и
второй системе шин. Порог срабатывания пуска по минимальному линейному напряжению задается уставкой
\textbf{«Uл КПОН 1(2)СШ»}, порог срабатывания пуска по максимальному напряжению обратной последовательности
задается уставкой \textbf{«U2 КПОН 1(2)СШ»}. Ввод в работу КПОН 1(2)СШ осуществляется с помощью программного
ключа \textbf{«КПОН 1(2)СШ»}.  

\item
Выбор пуска по напряжению КПОН осуществляется с помощью программного переключателя \textbf{«Тип КПОН 1(2)СШ»}: 
\begin{itemize}
\item \textit{«Umin, U2»} -- Пуск по минимальному линейному напряжению или максимальному напряжению обратной последовательности;
\item \textit{«Umin»} -- Пуск по минимальному линейному напряжению.
\end{itemize}

\item
Предусмотрена возможность контроля исправности цепей напряжения от внутренней логики БНН. При
появлении неисправности в цепях напряжения, действие логики КПОН задается программным
переключателем \textbf{«Неиспр ТН 1(2)СШ»}:
\begin{itemize}
\item \textit{«Выведено»} -- Неисправность в цепях ТН не вызывает изменения логики КПОН;
\item \textit{«Пуск КПОН»} -- При появлении неисправности в цепях ТН, КПОН всегда разрешен;
\item \textit{«Блок КПОН»} -- При появлении неисправности в цепях ТН, КПОН блокируется до исчезновения неисправности ТН.
\end{itemize}

\item
Оперативный ввод/вывод КПОН 1(2)СШ производится с помощью БУ (блоков управления функциями) \textit{«Вывод КПОН 1(2)СШ»}. 

\small
\begin{longtable}{|>{\centering\arraybackslash}m{5.3cm}|>{\centering\arraybackslash}m{4.2cm}|>{\centering\arraybackslash}m{1cm}|>{\centering\arraybackslash}m{3.3cm}|>{\centering\arraybackslash}m{1.8cm}|}
% Первая страница: полный заголовок
\caption{Перечень уставочных параметров КПОН\hfill\vspace{-0.5\baselineskip}}\label{kpon:tbl1}\\
\hline
\rowcolor{gray!30}
Наименование уставки & Диапазон & Шаг & Значение по умолчанию & Единица измерения \\
\hline
\endfirsthead

% Промежуточные страницы: «Продолжение таблицы»
\caption*{\hspace{3pt}\emph{Продолжение таблицы \ref{kpon:tbl1}}\hfill\ }\vspace{-0.5\baselineskip} \\
\hline
\rowcolor{gray!30}
Наименование уставки & Диапазон & Шаг & Значение по умолчанию & Единица измерения \\
\hline
\endhead

%===m>КПОН|-
%===m
\end{longtable}
\normalsize	

\end{enumerate}